% ECONOMICS_PROBLEM: {"id": "J62-497", "title": "Social connections as a causal mobility tool", "jel_code": "J62", "source_id": "OP-0366"}
% !TeX program = xelatex
\documentclass[11pt,a4paper]{article}
\usepackage[margin=23mm]{geometry}
\usepackage{fontspec}
\setmainfont{Latin Modern Roman}
\usepackage{amsmath,amssymb,mathtools}
\usepackage{xcolor,enumitem,booktabs,longtable,array,xurl,hyperref}
\definecolor{muted}{HTML}{53636C}
\hypersetup{colorlinks=true,urlcolor=blue,linkcolor=blue}
\setlength{\parindent}{0pt}
\setlength{\parskip}{5pt}
\newcommand{\R}{\mathbb R}
\newcommand{\E}{\mathbb E}
\newcommand{\N}{\mathbb N}
\newcommand{\ind}{\mathbf 1}
\newcommand{\Var}{\operatorname{Var}}
\newcommand{\Cov}{\operatorname{Cov}}
\newcommand{\supp}{\operatorname{supp}}
\DeclareMathOperator*{\argmin}{arg\,min}
\DeclareMathOperator*{\argmax}{arg\,max}
\newcommand{\entrymeta}[3]{{\sffamily\small\color{muted}\textbf{\texttt{#1}}\quad|\quad #2\par #3\par}}
\newcommand{\Problemref}[1]{\texttt{#1}}
\newcommand{\JELref}[2]{\texttt{#2} (JEL \texttt{#1})}
\begin{document}
\section*{Social connections as a causal mobility tool}
\textbf{Problem ID:} \texttt{J62-497}\quad \textbf{JEL:} \texttt{J62}\par
% BEGIN ECONOMICS STATEMENT
\label{problem:OP-0366}
{\sffamily\small\color{muted}Original subject group: Inequality and economic mobility\par}
\label{definitions:problem:30007648}
\textbf{Audit classification:} Published research agenda. Full published author copy and all authors checked. Direct connectedness-intervention tests are future work; association alone does not identify friendship effects.
\subsection*{Definitions and precise research target}
\textbf{Complete editorial problem.} Fix a cohort of children from low-income families and a network-forming intervention \(z\in\{0,1\}\). The intervention is precisely specified, such as a common activity with randomized cross-income participation; it does not directly give children higher income. Let \(A_i(z)\) be the share of child \(i\)'s friends from a predeclared higher socioeconomic group, and let \(Y_i(z)\) be adult income rank relative to a fixed reference distribution. Define the total policy effect \(\tau=E[Y_i(1)-Y_i(0)]\) and the network first stage \(\phi=E[A_i(1)-A_i(0)]\). Estimate both and determine whether increased connectedness improves mobility. When \(\phi\ne0\), interpreting \(\tau/\phi\) as an effect of the continuous friendship share requires explicit exclusion, interference restrictions, and either a justified local-average-response model or a linear response restriction; exposure or program participation may affect outcomes through other routes. Measure candidate mechanisms including information, aspirations, mentorship, and referrals with separately justified designs. Account for friendship selection, persistence of connections, attrition, and displacement of opportunities for untreated children. The task adapts the authors' explicit call for direct tests of connectedness interventions. A strong area-level correlation, including correlation with causal place effects, does not by itself identify the causal effect of creating these friendships.

\textbf{Definitions and admissible scope.} Let \(G(z)\) be an undirected friendship graph on a baseline cohort, and define low and high socioeconomic groups from prereform parental status. For node \(i\) with positive degree, \(A_i(z)=\sum_jG_{ij}(z)1\{j\in H\}/\sum_jG_{ij}(z)\); for isolates, report an indicator and use a declared value rather than dropping them. Adult rank is computed from raw income using the fixed reference population and midrank rule below. A policy effect includes altered exposure and friendships; it does not mechanically equal an effect of one added edge. Group randomization and exposure mappings must specify interference because changing one child's friendships changes others' networks. Exclusion for an instrumental interpretation must exclude all non-network routes of program assignment. In the rank formula let \(y_i(z)\) denote raw adult income and define \(Y_i(z)=F_{\rm ref}(y_i(z)-)+\tfrac12\Pr_{\rm ref}(y=y_i(z))\), a fixed-reference midrank with an explicit tie convention. Use \(A_i=0\) for isolates and report their frequency separately. Baseline low-income group membership, graph node set, follow-up age, and assignment units are fixed. Potential outcomes include the full graph assignment, with \(z\) an abbreviation for the declared joint regime. The finite-horizon total effect \(\tau\) and first stage \(\phi\) are distinct identified quantities under that regime; the ratio has no universal per-friend causal interpretation.
\subsection*{Variants and assumption changes}
\begin{enumerate}
\item Direct intervention test (source-motivated): vary cross-income activities while measuring actual friendship formation and adult mobility, with a long follow-up and baseline-defined cohorts.
\item Scale and opportunity displacement (editorial): vary school or neighborhood saturation and measure outcomes of untreated peers. Compare effects per induced durable connection when referral or school opportunities are scarce; permit effects to depend on network position.
\end{enumerate}
\subsection*{References and source audit}
\textbf{Support and access.} Full published author copy and all authors checked. Direct connectedness-intervention tests are future work; association alone does not identify friendship effects. \textbf{Locator:} Discussion, Nature p. 120, final research paragraph.
\begin{enumerate}\raggedright
\item\hypertarget{definitions:source:30007648:1}{} Raj Chetty; Matthew O. Jackson; Theresa Kuchler; Johannes Stroebel; Nathaniel Hendren; Robert B. Fluegge; Sara Gong; Federico Gonzalez; Armelle Grondin; Matthew Jacob; Drew Johnston; Martin Koenen; Eduardo Laguna-Muggenburg; Florian Mudekereza; Tom Rutter; Nicolaj Thor; Wilbur Townsend; Ruby Zhang; Mike Bailey; Pablo Barberá; Monica Bhole; Nils Wernerfelt. 2022. \emph{Social capital I: measurement and associations with economic mobility}. Discussion, Nature printed p. 120 (PDF p. 13), final research-agenda paragraph.
\url{https://tomrutter42.github.io/assets/pdf/sc1.pdf}.
DOI: \url{https://doi.org/10.1038/s41586-022-04996-4}.
\end{enumerate}

\subsection*{Common conventions: Source mathematical conventions}

These conventions apply to each entry unless explicitly replaced.

\textbf{Models and identification.} A primitive model \(m\) specifies all probability laws, feasible choices, information, equilibrium and measurement rules used by its target. The admissible class \(\mathcal M\) is fixed independently of outcomes. If \(P_O(m)\) is its observable probability law and \(\tau(m)\) the target, its identified set at observable law \(P\) is
\[
 \mathcal I_\tau(P)=\{\tau(m):m\in\mathcal M,\ P_O(m)=P\}.
\]
Point identification means that this set is a singleton. An empty set signals incompatibility of data and assumptions. Coordinate bounds need not describe the joint set. Bounds are sharp only if every retained target is attainable or arbitrarily approachable by compatible models. Finite bounds require explicit restrictions on unobserved quantities; unrestricted confounding, hidden wealth, extrapolation or arbitrary equilibrium selection can yield uninformative sets.

\textbf{Causal interventions.} A potential outcome \(Y_i(p)\) is the outcome under a complete policy regime \(p\), including timing, financing, information and market spillovers. The target population and units are fixed before treatment. With interference, use the complete assignment vector or a justified exposure mapping; individual no-interference notation alone cannot identify an economy-wide intervention. A difference of mean potential outcomes does not require the cross-world joint distribution. Conditioning on post-treatment migration, survival, enrollment or employment changes the estimand and needs separate justification.

\textbf{Statistical statements.} An estimator, confidence set, forecast or test specifies a sampling law, model class and loss/coverage criterion. Uniform \((1-\alpha)\)-coverage means \(\inf_{m\in\mathcal M}\Pr_m\{\tau(m)\in C_n\}\geq1-\alpha\), or an explicitly asymptotic version. The whole parameter space trivially has coverage, so informativeness requires a stated width, risk or power objective. A forecasting improvement is relative to a named benchmark, information set and evaluation population.

\textbf{Equilibrium and optimization.} An equilibrium comprises feasible plans, optimality, beliefs where needed, budget constraints and all market/resource clearing. The primitives or a stated benchmark define each correspondence \(\mathcal E(p;m)\). Nonemptiness is assumed or proved, never inferred just from compact policy sets. A selected-equilibrium target uses a declared selection rule; a robust target quantifies over all permitted equilibria. Use suprema or infima unless attainment is established. An \(\varepsilon\)-optimal policy is feasible and within \(\varepsilon\) of the finite optimum value. Report an empty feasible set separately.

\textbf{Welfare and accounting.} Normative utility, interpersonal normalization, social weights, numeraire, discounting and the use of public receipts are primitives. Behavioral utility need not coincide with the welfare criterion. Household welfare includes ownership income and rents once; adding profits again to compensating variation generally double-counts them. A monetary equivalent is defined by an explicit transfer and price convention, with monotonicity and range conditions. Average effects, mechanism contributions, transfers and welfare losses are different targets. Infinite sums require convergence and dynamic feasible plans require terminal or no-Ponzi conditions.

\textbf{Source versions.} A working-paper date, revision date and journal publication year are distinguished. A DOI for a journal version is not automatically the DOI of an earlier manuscript. Locators refer to the stated version and printed pages where specified. A metadata-only check does not verify a claim about a full-text conclusion. Every entry's source subsection narrows the provenance claim accordingly.
% END ECONOMICS STATEMENT
\end{document}
